\begin{abstract}
Hybrid offline-to-online reinforcement learning for transit control is attractive because operators often have rich historical data but only imperfect simulators.
However, a bus holding policy tuned in one city may fail in another city unless the target simulator is manually calibrated, and such calibration is expensive.
We study this cross-city transfer problem for bus holding using four open transit city bundles: Singapore, Austin/CapMetro, Halifax Transit, and MBTA.
We propose Causal-Factored Cross-City Model Transfer (\ours), a mechanism-wise residual world-model adapter that decomposes dense simulator corrections into headway, demand/dwell, speed, and reward mechanisms, and selects source corrections using only unlabeled target-city simulator summaries.
Against a dense \Htwo-style residual baseline, \ours is evaluated under a strict leave-one-city-out protocol in which each of the four cities is the target exactly once.
On static-derived cross-city dynamics validation, unweighted \ours beats \Htwo on all four targets with mean total-MSE ratio $0.759$, while similarity-weighted \ours also wins on all four targets and lowers the mean ratio to $0.578$.
Single-source zero-shot transfer is intentionally harder and mixed, with 6/12 wins and mean ratio $0.927$, while few-shot target offline adaptation keeps the mechanism-factored model below the dense adapted baseline on most target-route budgets; the strongest few-shot variant is a low-dimensional target context-bias adapter rather than full mechanism refitting.
In the full-module validation, a source-only feature-set selector plus guarded source-city ensemble lowers mean total MSE from $884.9$ for dense \Htwo-style transfer and $851.9$ for rigid \ours to $836.1$ over six configured cross-city splits, while preserving $6/6$ wins against \Htwo; the guard uses unlabeled target summaries but no target transition labels.
A same-information dense \Htwo source ensemble improves to $846.5$ but remains weaker than the guarded mechanism-factored ensemble.
In one-step policy-proxy evaluation, \ours improves reward/regret on three of four targets and lowers bunching rate on all four targets.
Target-information, per-mechanism, ablation, route-bootstrap, source-subset, and generator-bias checks support the same main trend while showing that the advantage depends on multi-source diversity and should not be overstated as superiority on every small city subset.
On two public Austin raw APC/AVL slices, source-only \ours is more stable than dense residual transfer but still worse than a state-informed passive no-correction baseline; with few-shot target residual calibration, a residual gate beats that passive baseline at $1\%$ target line-group budget and at chronological budgets of six hours or more in both slices.
Because passive bus logs cannot identify unobserved holding-action counterfactuals, we add a counterfactual traffic-signal-control stress test on RESCO.
There, pressure-guarded \ours reaches MaxPressure-level mean queue on the 600-second multi-seed stress test and improves over dense residual transfer, simulator-only MPC, native actuated SUMO, and fixed signal programs, while a 3600-second check correctly exposes phase pressure as the stronger long-horizon classical rule.
Full-budget official RESCO and patched LibSignal RL baselines are executed as native-protocol leaderboard cross-checks, not as same-protocol rows for the \ours phase-MPC table.
The resulting contribution is a reproducible, reviewer-facing protocol and an empirically supported boundary condition: mechanism-factored transfer is useful for target-label-efficient cross-city residual adaptation, but counterfactual control or field validation remains necessary before operational claims.
\end{abstract}
